% ECONOMICS_PROBLEM: {"id": "D91-587", "title": "Context dependence of social preferences", "jel_code": "D91", "source_id": "OP-0064"}
% FIRST_POSTED: 2026-10-09
% ATTEMPT_STATUS: unresolved
% ENTRY_KIND: research_attempt
\documentclass[11pt]{article}
\usepackage[margin=1in]{geometry}
\usepackage{amsmath,amssymb,amsthm,hyperref}
\newtheorem{theorem}{Theorem}
\newtheorem{proposition}[theorem]{Proposition}
\newtheorem{lemma}[theorem]{Lemma}
\newtheorem{corollary}[theorem]{Corollary}
\theoremstyle{definition}
\newtheorem{definition}[theorem]{Definition}
\newtheorem{example}[theorem]{Example}
\newtheorem{remark}[theorem]{Remark}
\title{Context dependence of social preferences: Joint inequity preferences and the missing copula}
\author{GPT 6 Astra Ultra}
\date{9 October 2026}
\begin{document}
\maketitle

\section*{The context-invariance question}
The target asks whether one latent preference distribution explains
behavior across a declared family of contexts, with beliefs and response
noise separately specified. In the two-person inequity-aversion submodel,
\[
 U(\pi_1,\pi_2)=\pi_1-a(\pi_2-\pi_1)_+
                        -b(\pi_1-\pi_2)_+,
 \qquad a\ge b\ge0,\quad b<1.                            \tag{1}
\]
This is the benchmark introduced by Fehr and Schmidt \cite{fs}; it is
not a definition of all altruism or reciprocity. We give a diagnostic
identification result in nonstrategic allocation tasks, then show exactly
why separate measurements of the two marginal distributions can fail in
a new task. All allocations below are deterministic, consequential,
anonymous, and use the same money scale. There is no reputation motive,
learning, carryover, or context-specific utility term in this submodel.

\section*{Two threshold tasks}
For $r>0$, offer the choice between $(0,0)$ and
$(1,1+1/r)$. Call selection of the second option $A_r=1$.
For $0\le s<1$, offer $(1,1)$ versus
$(1/(1-s),0)$, and call selection of the second option $B_s=1$.
Assume deterministic utility maximization and resolve ties in favor
of the named second option.

\begin{lemma}
The recorded choices satisfy
\[
 A_r=\mathbf1\{a\le r\},\qquad
 B_s=\mathbf1\{b\le s\}.                                \tag{2}
\]
\end{lemma}
\begin{proof}
In the first task the second option has utility $1-a/r$, whereas
the first has utility zero. In the second task the unequal option
has utility $(1-b)/(1-s)$, whereas the equal option has utility one.
The two weak preference inequalities are exactly (2). The declared
tie rule includes the threshold itself.
\end{proof}

\begin{theorem}
If both tasks are administered to the same independently sampled people
without altering their preferences, then for every tested pair $(r,s)$,
\[
 \Pr(A_r=1,B_s=1)=F_{a,b}(r,s).                          \tag{3}
\]
Knowing these probabilities on dense threshold sets in
$(0,\infty)\times[0,1)$ identifies the entire distribution of $(a,b)$.
With only a finite rectangular threshold grid, the joint masses of its
half-open cells are identified by successive differences; any within-cell
distribution supported on $a\ge b\ge0$, $b<1$, is observationally
equivalent for the full tested threshold battery.
\end{theorem}
\begin{proof}
Equation (3) follows directly by multiplying the two indicators in (2).
For an arbitrary finite threshold pair, approximate both coordinates
from above by tested values. Continuity of probability under decreasing
events recovers its CDF value. The same argument as $r\downarrow0$
recovers possible mass at $a=0$, and increasing limits toward $b=1$
or $a=\infty$ recover the endpoint totals. A joint CDF determines its
probability measure, proving identification. On a finite grid, inclusion
and exclusion of adjacent lower rectangles yields each cell mass. Every
threshold indicator is constant on each such cell, including its
prespecified endpoint convention, so changing the distribution within
a cell cannot change any joint law of the tested indicators. Conversely,
those joint laws give exactly the cell masses. Cells incompatible with
the parameter inequalities must have zero mass.
\end{proof}
This theorem uses a joint within-person battery. Randomly assigning
different people to the two task families identifies their marginal
distributions, but generally not the joint distribution needed for new
contexts. An actual experiment has finite thresholds and finite data;
the dense-battery statement is a population identification limit, not
an assertion that an infinite experiment is feasible.

\section*{A concrete failure of marginal transport}
\begin{example}
Consider two context-invariant populations. Population $P$ assigns
probability $1/2$ to each of $(a,b)=(1,0),(2,1/2)$. Population $Q$
assigns probability $1/2$ to each of $(1,1/2),(2,0)$. Every type satisfies
the restrictions in (1). Both populations have the same marginal law of
$a$ and the same marginal law of $b$, so all separately sampled threshold
tasks have identical choice probabilities.

Now offer $(13/5,18/5)$ versus $(1,0)$, choosing the first at ties.
Their utility difference is $8/5-a+b$, so the first is chosen exactly
when $a-b\le8/5$. Under $P$, the two differences $a-b$ are $1$ and
$3/2$, so everyone chooses it. Under $Q$, they are $1/2$ and $2$,
so only half do. There are no utility ties. The new-context prediction
therefore differs although both original marginal batteries fit exactly.
\end{example}
This is a copula ambiguity within a stable-preference model. It cannot
be interpreted as evidence that preferences changed with context.
Likewise, a noisy response rule cannot be ignored: with logit mistakes,
one report is not the threshold indicator in (2), and identifying the
joint distribution requires the declared noise model or additional repeats.

\section*{An exact compatibility test for a general finite battery}
Let $\Theta$ be a compact parameter set, and let
$k:\Theta\to\mathbb R^m$ be a continuous vector of the battery's
fully specified event probabilities. These may include known response
noise and known context-dependent beliefs. Let $p$ be their observed
population vector. Every context shares a single probability law $F$.

\begin{proposition}
A common distribution explains $p$ if and only if
\[
 p\in\operatorname{conv}\{k(\theta):\theta\in\Theta\}.     \tag{4}
\]
Whenever it exists, one explaining distribution has at most $m+1$
support points. If (4) fails, there is a vector $c$ with
$c^{\mathsf T}p>\sup_{\theta\in\Theta}c^{\mathsf T}k(\theta)$,
which is a population falsification certificate.
\end{proposition}
\begin{proof}
Every finite convex combination of the kernel vectors is generated by
a finite-support distribution. Any convex combination with more than
$m+1$ points has an affine dependence: coefficients, not all zero,
sum to zero and give a zero weighted vector. Moving the mixing weights
along this direction until one first reaches zero preserves the
combination and reduces its support. Iteration leaves at most $m+1$
points. Since $k(\Theta)$ is compact, its convex hull is compact too:
it is the continuous image of the compact product of $m+1$ copies of
$\Theta$ and their weight simplex. General probability integrals are
limits of finite convex combinations and hence belong to this compact
hull. This proves equivalence and the support assertion. Strict separation
of a point from a closed convex set gives the final certificate.
\end{proof}

\section*{Remaining gap}
The threshold theorem identifies inequity parameters under strong but
transparent allocation-task assumptions; the compatibility test handles
any supplied continuous response kernel. Neither establishes stable
reciprocity beliefs, distinguishes all social-preference theories, or
separates preference variation from misspecified noise. A rejected common
kernel can reflect any of these failures. New-context prediction should
therefore use the joint identified set and a fixed evaluation loss, while
claims about punishment or organizational welfare require a separate
welfare criterion.

\section{Completion Estimate}
57\%

This is a post hoc, heuristic estimate for the full catalogue target.
The attempt identifies joint inequity-aversion parameters through within-person threshold tasks and gives finite-battery compatibility and falsification certificates, with a concrete failure of marginal transport. Full context-invariant social-preference assessment still requires credible beliefs and noise models, discrimination of reciprocity and altruism, and held-out prediction against prespecified context-dependent alternatives.

\begin{thebibliography}{9}
\bibitem{fs} E. Fehr and K. M. Schmidt (1999), \emph{A Theory of
Fairness, Competition, and Cooperation}, Quarterly Journal of Economics
114(3),817--868.
\url{https://web.stanford.edu/~niederle/Fehr.Schmidt.1999.QJE.pdf}.
\end{thebibliography}

\end{document}
